\documentclass[10pt,a4paper]{amsart}
\usepackage[margin=19mm]{geometry}
\usepackage{amsmath,amssymb,mathtools}
\usepackage[scr=rsfs,cal=euler]{mathalfa}
\usepackage{xfrac}
\usepackage[unicode,hidelinks,bookmarks=false]{hyperref}
\newcommand{\ie}{i.e.\ }
\newcommand{\cf}{cf.\ }
\newcommand{\resp}{respectively\ }
\newcommand{\Subsection}{\subsection}
\providecommand{\nobreakdash}{\nobreak\hbox{-}}
\setlength{\emergencystretch}{2em}
\multlinegap0pt
\usepackage{bm}
\usepackage{bbm}
\usepackage{dsfont}
\usepackage{xspace}
\usepackage{cancel}
\usepackage{mathtools}
\usepackage{amsthm}
\makeatletter
\@ifundefined{theorem}{\newtheorem{theorem}{Theorem}}{}
\makeatother
\title[On Proximity and other Distance Parameters in Planar Graphs]{On Proximity and other Distance Parameters in Planar Graphs}
\author{Peter Dankelmann, Sonwabile Mafunda, Sufiyan Mallu}
\date{Source: arXiv:2508.10078v1 (CC BY 4.0)}
\begin{document}
\maketitle

\noindent\emph{Source and licence.} On Proximity and other Distance Parameters in Planar Graphs. Version arXiv:2508.10078v1 (CC BY 4.0), distributed under the Creative Commons Attribution 4.0 International licence, as recorded in the arXiv licence metadata for this exact version and shown on the abstract page https://arxiv.org/abs/2508.10078v1. Full rights notice: licenses/arxiv-2508.10078-CC-BY-4.0.txt.\par\smallskip

\noindent\textit{Editorial scope.} Counted unit: the Theorem whose source label is \texttt{theo:remoteness-proximity} in the article (source0.tex lines 1011--1017, source label \texttt{theo:remoteness-proximity}), delivered together with the article's own complete original proof of it (source0.tex lines 1020--1108, one \texttt{proof} environment of 89 lines). No mathematics is written, repaired or re-derived: statement and proof are the article's own text line for line. Required context retained uncounted: none. Every definition, display and label of those lines is retained unchanged. Omitted from this package and not redistributed: 1--1010, 1018--1019, 1109--1638. Nothing inside a retained excerpt is omitted or altered: every retained body is delivered exactly as the article's lines are written, so no omission receipt of any kind accompanies this package. Citations: the retained excerpts carry no citation command at all, so no bibliography entry of the article is transcribed and no reference is redistributed. Reference adaptation: none was needed and none was made; every reference inside the delivered unit resolves inside it, so no unresolved reference or citation is produced. Printed slips kept literal: no printed word, symbol, hypothesis or number of the counted statement or of its proof was altered. Layout and class: the article's own class is \texttt{article} with packages including authblk, babel, amsmath, amsfonts, amsthm, amssymb, pgf, tikz, pgfplots, float, graphicx, tikz-cd; the wrapper is the lane's pinned \texttt{amsart} editorial wrapper at 10pt on A4, which restates the theorem-like environments the retained text needs and adds the article's own macro definitions that the retained text needs (none), with extra wrapper packages (bm, bbm, dsfont, xspace, cancel, mathtools). The delivered numbering is the unit's own. Assets: no retained span contains a figure environment, a graphics-inclusion command, a PDF-page inclusion, a table or a TikZ picture; no image file is copied into this package, the package declares no asset dependency and no image file is redistributed with it. Licence: version arXiv:2508.10078v1 of this article is distributed under the Creative Commons Attribution 4.0 International licence, as recorded in the arXiv licence metadata for this exact version and shown on the abstract page https://arxiv.org/abs/2508.10078v1. A full rights notice is delivered at licenses/arxiv-2508.10078-CC-BY-4.0.txt. Characters: every retained excerpt is pure ASCII, so the delivered engine, XeLaTeX with its default Latin Modern text font, reports no missing character.
\begin{theorem}\label{theo:remoteness-proximity}
Let $n, \kappa \in \mathbb{N}$ with $\kappa \leq \frac{n+1}{2}$. Let $G$ be a graph of order 
$n$ with $\kappa(G) = \kappa$. Then
\[
\rho(G) - \pi(G) \leq \frac{n+2\kappa-3}{4\kappa} 
              - \frac{(3\kappa+1)(\kappa-1)}{4\kappa (n-1)}.       \]                           
\end{theorem}

\begin{proof}  
Let \( u \) and \( v \) be vertices of \( G \) such that \( \overline{\sigma}(u) = \rho(G) \) 
and \( \overline{\sigma}(v) = \pi(G) \), and let $k:= d(u,v)$. For $i \in \mathbb{N} \cup \{0\}$ let 
\( N_i(u) \) denote the set of vertices at distance \( i \) from \( u \). 
Our first objective is to find a set \( X \subseteq V(G) \) such that \( \sigma(u|X) - \sigma(v|X) \) 
is small. We then bound \( \sigma(u|V(G) - X) - \sigma(v|V(G) - X) \).

Since $G$ is $\kappa$-connected, we have $|N_i(u)| \geq \kappa$ for all 
$i\in \{1,2,\ldots,{\rm ecc}(u)-1\}$. For \( i \in \{1, 2, \ldots, k-1\} \), define $M_i$ 
to be a set of exactly $\kappa$ vertices of $N_i(u)$, and let \( M_0 =\{u\} \) and \(M_k = \{ v\} \).

Clearly, $M_i \cap M_j = \emptyset$ for all $i, j \in \{0, 1,\ldots,k\}$ with $i\neq j$. 
For \( i \in \{0, 1, \ldots, \lfloor\frac{k}{2}\rfloor\} \) define the sets 
\( B_i = M_i \cup M_{k-i} \) . Let 
\( B = \bigcup_{i=0}^{\lfloor\frac{k}{2}\rfloor} B_i \), so \( |B| = \kappa(k - 1)+2 \).

Note that for each \( y \in M_i \), we have \( d(u, y) = i \) and \( k-i \leq d(v, y) \), and for each \( y \in M_{k-i} \), we have \( d(u, y) = k - i \) and \( i \leq d(v, y)  \). Hence, for $i \neq 0,k$, 
\begin{eqnarray*}
\sigma(u|B_i) - \sigma(v|B_i) 
  &=&  \sigma(u|M_i) + \sigma(u|M_{k-i}) - \sigma(v|M_i) - \sigma(v|M_{k-i}) \\
& \leq&  \kappa i + \kappa(k-i) - \kappa(k-i) - \kappa i \\
& =&  0.
\end{eqnarray*}
By summing the above equation for \( i \in \{1, \ldots, \lfloor\frac{k}{2}\rfloor\} \), and noting 
that \(\sigma(u|B_0) = d(u,v) = \sigma(v|B_0) \), we obtain that 
\begin{equation}\label{eq:total-distance-of-B}
\sigma(u|B) - \sigma(v|B) \leq 0.
\end{equation}
Let $B'=V(G)-B$. Now consider \( w \in B' \). We claim that
\begin{equation} \label{eq:vertices-notin-B} 
d(u,w) - d(v,w) \leq \left\{ \begin{array}{cl}
k-1 & \textrm{if $w \in B' \cap N_{\leq k}(u)$,} \\
k & \textrm{if $w \in B' \cap N_{\geq k+1}(u)$.}
\end{array} \right. 
\end{equation}
Indeed, if $w \in B' \cap N_{\leq k}(u)$, then $d(u,w) \leq k$ and $d(v,w) \geq 1$, so 
$d(u,w) - d(v,w) \leq k-1$. If $w \in B' \cap N_{\geq k+1}(u)$, then 
by the triangle inequality, we have
\[
d(u, w) \leq d(u, v) + d(v, w),
\]
and \eqref{eq:vertices-notin-B} follows since $d(u,v)=k$. 


We first show that the theorem holds if $n \leq k\, \kappa$, that is, if there are only few vertices
besides those in $B$. In this case we have  
${\rm diam}(G) \leq 1 + \frac{n-2}{\kappa} \leq 1 + \frac{k\kappa-2}{\kappa} < 1+k$,
and so all vertices of $B'$ are in $B' \cap N_{\leq k}(u)$. Hence by 
\eqref{eq:total-distance-of-B} and \eqref{eq:vertices-notin-B} we obtain
\[ \sigma(u) - \sigma(v) = \sigma(u|B) - \sigma(v|B) + \sum_{w \in B'} \big( d(u,w)-d(v,w)\big)  
           \leq |B'|(k-1). \]
Now $|B'|=n-((k-1)\kappa+2) \leq \kappa-2$ and $n \geq (k-1)\kappa+2$, which implies that
$k-1 \leq \frac{n-2}{\kappa}$. Substituting these bounds into the above inequalities and 
dividing by $n-1$ implies that
\[ \rho(G) - \pi(G) \leq \frac{\kappa-2}{\kappa} \frac{n-2}{n-1}, \]
which by a simple calculation,
yields the theorem. Hence we may assume from now on that $n \geq k \kappa$. 

Since $n > k \kappa$, we have $|B'| \geq \kappa-1$.
We claim that the set $B'$ has at least $\kappa-1$ vertices in $N_{\leq k}(u)$. 
Indeed, if less than $\kappa-1$ vertices of $B'$ are in $N_{\leq k}(u)$, then 
$N_{\geq k+1}$ is nonempty, and the set $N_k(u)$ is a cut-set containing not more than $\kappa-2$ 
vertex besides $v$, a contradiction to $G$ being $\kappa$-connected. 
Summing \eqref{eq:vertices-notin-B} over all \( w \in B' \), we thus get
\begin{eqnarray} 
\sigma(u|B') - \sigma(v|B')  
 & \leq & \sum_{w \in B' \cap N_{\leq k}} (k-1) + \sum_{w \in B' \cap N_{\geq k+1}} k \nonumber \\
 & \leq & k(n-|B|)-(\kappa-1) \nonumber \\ 
 & = & k(n - (k-1)\kappa-2)-\kappa+1. \label{eq:total-distance-of-B'}
\end{eqnarray}
Combining equations (\ref{eq:total-distance-of-B}) and (\ref{eq:total-distance-of-B'}), we obtain that
\begin{eqnarray}
\sigma(u) - \sigma(v) 
& = & \sigma(u|B) - \sigma(v|B) +\sigma(u|B') - \sigma(v|B') \nonumber  \\
& \leq &   k(n - (k-1)\kappa-2)-\kappa+1. \label{eq:rho-pi-1}
\end{eqnarray}
A straightforward maximisation shows that for \( k = \frac{n+\kappa-2}{2\kappa} \), the right-hand side of the above 
inequality is maximised. Substituting this value of \( k \) into inequality 
\eqref{eq:rho-pi-1}, we obtain
\[
\sigma(u) - \sigma(v) \leq \Big( \frac{n+\kappa-2}{2\kappa} \Big)^2 - \kappa + 1.
\]
Dividing by $n-1$ in the above inequality yields, after simplification,
\[
\rho(G) - \pi(G) \leq \frac{n+2\kappa-3}{4\kappa}   
                 - \frac{(\kappa-1)(3\kappa+1)}{4\kappa (n-1)}.
\]
This completes the proof of the theorem.
\end{proof}

\end{document}
